% This text is a supplement of CINTA.
% May. 2018
% Author: Libin Wang
% SCNU
% additional use of \usepackage{beamerthemesplit}
\documentclass{beamer}

%%%%%%%%%%%%%%%%%%%%%%%%%%%%%%%%%%%%%%%%%%%%%%%%%%%%%%%%%%%%%%%%%%%%%
\usepackage{xeCJK}
%\usepackage[CJKbookmarks]{hyperref}

%% Support Chinese with TeX Live fonts.

\setCJKmainfont[
  BoldFont=FandolHei-Regular.otf,
  ItalicFont=FandolKai-Regular.otf,
  BoldItalicFont=FandolKai-Regular.otf
]{FandolSong-Regular.otf}
\setCJKsansfont[AutoFakeBold, AutoFakeSlant]{FandolHei-Regular.otf}
\setCJKmonofont{FandolFang-Regular.otf}

%\setCJKmainfont[BoldFont=Microsoft YaHei,ItalicFont=YouYuan]{SimSun} %设置默认中文字体

\setmainfont{Times New Roman} % 设置默认英文衬线字体

\setCJKfamilyfont{ls}{FandolSong-Regular.otf} %定义字体
%%%%%%%%%%%%%%%%%%%%%%%%%%%%%%%%%%%%%%%%%%%%%%%%%%%%%%%%%%%%%%%%%%%%%

%\usepackage{beamerthemesplit} % new 
\usepackage{beamerthemeshadow}
\setbeamertemplate{headline}{}
\setbeamertemplate{footline}{}
\setbeamertemplate{navigation symbols}{}
\usepackage{amsmath, amsfonts, amssymb}

%% Setup for Code
\usepackage{listings}
\lstset{language = Python,
	showstringspaces = false,
	columns = fullflexible,
	numbers = left,
	numberstyle = \tiny,
	frame = single}
%Can't work!
%\usepackage{minted}

%%\newproof{pf}{Proof} 
%\newtheorem{thm}{Theorem}[section]
%\newtheorem{fact}[theorem]{Fact}
\newtheorem{proposition}[theorem]{Proposition}
\newtheorem{remark}[theorem]{Remark}
%\title{A Concrete Introduction to Number Theory and Algebra--CRT} 
\title{Mathematics for Computer Science -- CRT} 
\author{Libin Wang} 
\institute{School of Computer Science, South China Normal University}
\date{\today} 

\begin{document}

\frame{\titlepage} 

\frame{\frametitle{Table of contents}\tableofcontents} 


\section{The Chinese Remainder Theorem（中国剩余定理）-- 数论视角.} 
\frame{\frametitle{导引.} 
数论部分，主要目标是解方程。群论部分，主要目标是群结构。

\begin{block}{思路.}
\begin{itemize}
	\item 数论部分，目前我们学习解方程方法包括：一元一次同余方程，一元$n$次同余方程。接下来需要扩展到一元同余方程组的解法。未来是多元同余方程等等。
	\item 群论部分，主要目标是群结构。群结构与解同余方程有什么关系吗？
\end{itemize}
\end{block}
\pause

\begin{alertblock}{目标}
	给出一种求解一元一次同余方程组的高效算法；探讨群结构与解同余方程组的关系。
\end{alertblock}


}

\frame{\frametitle{导引.} 
Chinese Remainder Theorem（中国剩余定理），或称为\emph{中国余数定理}则更准确。

\begin{block}{历史.}
中国剩余定理有着满满的中国传统文化元素，与之相关的名人包括：春秋时期的孙子、秦汉时期的韩信、宋朝的沈括和秦久韶等。
南北朝时期（公元5世纪）的数学著作《孙子算经》卷下第二十六题所谓“物不知数”问题则为国人所熟知：
有物不知其数，三三数之剩二，五五数之剩三，七七数之剩二。问物几何？
\end{block}
}

\frame{\frametitle{Motivation.} 
中国剩余定理
讨论一元同余方程组的高效解法。
	\begin{example}[Example 1.]
	Suppose we wish to solve:
\begin{align*}
		    x \equiv  2\pmod{5}\\
		    x \equiv  3\pmod{7}
\end{align*}
	\end{example} \pause

\begin{block}{Solution.}
	1. We have $x = 5t + 2$ from the first congruence, for $t \in \mathbb{Z}$;\pause
	
	2. Substitute for $x$ in the second congruence, $ 5t + 2 \equiv 3  \pmod{7}$;\pause
	
	3. Simplifying, we get $5t \equiv 1 \pmod{7}$;\pause
	
	4. Multiplying both sides by $5^{-1}$, we get $t = 7s + 3$ for $s \in \mathbb{Z}$;\pause
	
	5. Finally, $x = 35s + 17$, which means $x \equiv 17 \pmod{35}$.
	
\end{block}
}

\frame{\frametitle{Question.} 
	\begin{block}{思考.}
	以上求解一元同余方程组的方法是否高效？如果不高效，能否给出更高效的算法？
	\end{block} 
}

\frame{\frametitle{The Chinese Remainder Theorem--CRT.} 
	For any system of equations like this, the \emph{Chinese Remainder Theorem}, short for CRT, tells us there is always a unique solution up to a certain modulus, and describes how to find the solution efficiently.
	\begin{theorem}
	     Let $p, q$ be primes, $n = pq$. For each $a \in \mathbb{Z}_p$, $b \in \mathbb{Z}_q$, there is a unique $x$, $0 \leq x < n$ such that
	     $x \equiv a\pmod{p}$ and $x \equiv b\pmod{q}$.
	\end{theorem}
}

\frame{\frametitle{The Chinese Remainder Theorem--CRT.} 
	\begin{theorem}
		Let $p, q$ be two relatively prime positive integers, $n = pq$. For each $a \in \mathbb{Z}_p$, $b \in \mathbb{Z}_q$, there is a unique $x$, $0 \leq x < n$ such that
		$x \equiv a\pmod{p}$ and $x \equiv b\pmod{q}$.
	\end{theorem}
Proof Idea
\begin{enumerate}
\item Given $a$ and $p$, how can we find some $c$ s.t. $a c \equiv a \pmod{p}$? \pause

\item $c$ must be some $1$ under modulo $p$\pause

\item 模$p$下为$1$的数可能有哪些（可参考CINTA 附录D.）? \pause

\item Find some $c$ s.t. $a c c^{-1} \equiv a \pmod{p}$ \pause

\item Then $x$ must be something like that $x = (a c c^{-1} + b d d^{-1})$, what should be $c$ and $d$?
\end{enumerate}
}

\frame{\frametitle{The Chinese Remainder Theorem--CRT.} 
	\begin{theorem}
		Let $p, q$ be coprime positive integers, $n = pq$. For each $a \in \mathbb{Z}_p$, $b \in \mathbb{Z}_q$, there is a unique $x$, $0 \leq x < n$ such that
		$x \equiv a \pmod{p}$ and $x \equiv b \pmod{q}$.
	\end{theorem}

\begin{proof}
	By construction. Since $p, q$ are coprime, there exist $p_1$ and $q_1$ such that $p_1 \equiv p^{-1}  \pmod{q}$ and $q_1 \equiv q^{-1} \pmod{p}$. Define $x$ by:
		\[
		x = a q q_1 + b p p_1
		\]
	It is easy to check that $x$ satisfies both equations.
	It remains to show no other solutions exist modulo $n$. Suppose $\exists z \neq x$ is another solution. Then $(z - x ) = tp$ and $(z - x) = sq$, for some $t, s \in \mathbb{Z}$. 
	Since $p$ and $q$ are coprime, then $(z - x) = kpq$, for $k \in \mathbb{Z}$. Hence $z \equiv x \pmod{n}$. 
	
\end{proof}
}

\frame{\frametitle{Example.} 
	\begin{example}[Example 2.]
		Suppose we wish to solve:
		
		\begin{align*}
		    x \equiv  2\pmod{5}\\
		    x \equiv  3\pmod{7}
		\end{align*}
		
	\end{example}

\begin{block}{Solution.}
	1. Let $a=2$, $b = 3$, $p = 5$, $q = 7$, $n = pq = 35$;
	
	2. Compute $p_1 \equiv p^{-1} \pmod{q}$ and $q_1 \equiv q^{-1}\pmod{p}$ using $\operatorname{egcd}$ algorithm; $p_1 = 3 $, $q_1 = 3$;
	
	3. $ x \equiv a q q_1 + b p p_1 \pmod{n}$; $ x \equiv 17 \pmod{35}$;
	
	4. It is easy to check that $x$ is a correct solution.	
\end{block}
}

\frame{\frametitle{Exercise.} 
	\begin{block}{求解以下方程.}
		Suppose we wish to solve:
		\begin{align*}
				    x \equiv  3\pmod{11}\\
				    x \equiv  4\pmod{13}
		\end{align*}
		
	\end{block}
\pause
\begin{block}{Ans.}
	$69$
\end{block}
}

\frame{\frametitle{Generalization.} 
	For several equations, we have a generalized version of CRT.
\begin{theorem}
	Let $m_1, m_2, \ldots, m_n$ be a set of pairwise relatively prime positive integers. Then the system of $n$ equations:
	\begin{align*}
	  x &\equiv a_1 \pmod{m_1}\\
	    &\cdots \\
	  x &\equiv a_n \pmod{m_n}
	\end{align*}
	has a unique solution for $x$ modulo $M$ where $M = m_1 m_2 \ldots m_n$.
\end{theorem}	

}

\frame{\frametitle{Generalization.} 
	\begin{proof}
		By construction. Let $M= \prod_{i = 1}^{n} m_i$, $b_i = M/m_i$, $b_i' = b_i ^{-1} \pmod{m_i}$. Then 
		\[
		y = \sum_{i=1}^{n} a_i b_i b_i' \pmod{M}
		\]
		is the unique solution.
	\end{proof}
}

\frame{\frametitle{思考题-1.}
\begin{block}{思考题.}
设$m$和$n$为互素的正整数，$a>0$为一个正整数，如果
\begin{align*}
x &\equiv a \pmod{m}\\
x &\equiv a \pmod{n}
\end{align*}
$x$模$mn$等于什么？为什么？
\end{block}
}

\frame{\frametitle{思考题-2.} 
	如果同余方程组的模数$m_1, m_2, \ldots, m_n$不两两互质，那么上述方程组是否还有解？如果有解，怎么求？请思考下面这题。
\begin{block}{思考题.}
	判定下列两个同余方程组是否有解。若有，求出模适当模数下的所有解。
\[
\text{(a)}\;
\begin{cases}
x \equiv 7 \pmod{12} \\
x \equiv 3 \pmod{8}
\end{cases}
\qquad\qquad
\text{(b)}\;
\begin{cases}
x \equiv 5 \pmod{12} \\
x \equiv 2 \pmod{8}
\end{cases}
\]
\end{block}	
\pause

\begin{block}{思路}
	模数不互素，也就是说有大于$1$的公共因子$d$，这个$d$能起到什么作用？其次，反向应用思考题-1的结论，从一个方程式化为多个方程式。
\end{block}
}

\section{CRT的抽象代数视角.}

\frame{\frametitle{A perspective from Abstract Algebra.} 
提醒：本节关注的不是算法，而是CRT背后的思想。
	\begin{block}{Motivation.}
		Let $n = pq$, where $p,q > 1$ are relatively prime. Given a positive integer $x$, it can be expressed as a unique pair $\bigl([x \bmod p], [x \bmod q ]\bigr)$.
	\end{block}
}

\frame{\frametitle{A perspective from Abstract Algebra.} 
\begin{theorem}
	Let $p, q > 1$ be coprime, $n = pq$. Then
	\[
	\mathbb{Z}_n \cong \mathbb{Z}_p \times \mathbb{Z}_q \quad (\text{as an additive group})\] 
	and,
	\[ \quad \mathbb{Z}_n^* \cong \mathbb{Z}_p^* \times \mathbb{Z}_q^* \quad (\text{as a multiplicative group}).
	\]
\end{theorem}

\begin{proof}
1. Define $\phi: \mathbb{Z}_n \to \mathbb{Z}_p \times \mathbb{Z}_q$ by:
\[
\phi(x) \triangleq \bigl([x \bmod p ], [x \bmod q]\bigr)
\]

2. Show $\phi$ is bijective.

3. Check that $\phi(x)$ preserves the group operation.\pause

The proof that it is an isomorphism from $\mathbb{Z}_n^*$ to $\mathbb{Z}_p^* \times \mathbb{Z}_q^*$ is similar.
\end{proof}
}

\frame{\frametitle{Remark.}
\begin{remark}
首先，注意到$\mathbb{Z}_p \times \mathbb{Z}_q$是加法群，请验证！
\end{remark}
}

\frame{\frametitle{How to prove the bijection?}
\begin{proof}
证明映射$\phi$是一种双射，即证明$\phi$是满射且单射。满射显然，因为根据
中国剩余定理，任意序对中的两个同余式在模$n$下存在唯一解。证明单射即证明，
如果对任意整数$0 \le a, b < n$，有
$\bigl([a \bmod p ], [a \bmod q ]\bigr) = \bigl([b \bmod p], [b \bmod q]\bigr)$，则$a = b$。再次根据
中国剩余定理可得。
\end{proof}
} 

\frame{\frametitle{How to prove the isomorphism?}
\begin{proof}
证明映射$\phi$保持群操作，即需要证明：
\begin{align*}
\phi(a+b) &= \bigl([(a+b) \bmod p ], [(a+b) \bmod q]\bigr) \\
          &= \bigl([a \bmod p ], [a \bmod q] \bigr) + \bigl([b \bmod p ], [b \bmod q ]\bigr)\\
          &= \phi(a) + \phi(b)
\end{align*}
\end{proof}
}

\frame{\frametitle{Example.} 
	\begin{example}[Example 3.]
		Take $n = 15 = 5 \cdot 3$.
		$\mathbb{Z}_n^* = \{ 1, 2, 4, 7, 8, 11, 13, 14\}$ is isomorphic to $\mathbb{Z}_5^* \times \mathbb{Z}_3^*$ since we can give the following correspondence:
		\[
		1 \leftrightarrow (1, 1) \quad 2 \leftrightarrow(2,2) \quad 4 \leftrightarrow (4,1) \quad 7 \leftrightarrow (2,1)
		\]
		\[
		8 \leftrightarrow (3,2) \quad 11\leftrightarrow(1,2) \quad 13 \leftrightarrow (3,1) \quad 14 \leftrightarrow (4,2)
		\]
	\end{example}
}

\frame{\frametitle{Example.} 
	\begin{example}[Example 4.]
		To compute $14\cdot 13 \bmod 15$. Since $14 \leftrightarrow (4, 2)$ and $13 \leftrightarrow (3, 1)$, we have:
		\[
		(4, 2) \cdot (3, 1) = ([4\cdot 3 \bmod 5], [2\cdot 1\bmod 3]) = (2,2).
		\]
		Note that $(2, 2) \leftrightarrow 2$, which is the correct answer.
	\end{example}
}

\frame{\frametitle{Example.} 
	\begin{example}[Example 5.]
		To compute $11^{53} \bmod 15$. Since $11 \leftrightarrow (1, 2)$ and $2 \equiv -1  \bmod 3$ we have:
		\[
		(1, 2)^{53} = ([1^{53} \bmod 5], [(-1)^{53} \bmod 3]) = (1, -1 \bmod 3) = (1, 2).
		\]
	 Thus, $11^{53} \bmod 15 = 11$ 
	\end{example}
}

\frame{\frametitle{Example.} 
\begin{example}[Example 6.]
\label{exm:eq_sqr_cong}
	设$n = p \cdot q$为合数，$p$和$q$是两个不同的奇素数。以下等式在模$n$的意义上有多少个解？
	\[ x^2 \equiv 1 \pmod{n}\]
为此，需要解以下两个同余方程
	\begin{align*}
	x^2 &\equiv 1 \pmod{p} \\
	x^2 &\equiv 1 \pmod{q}
	\end{align*}
以上两式分别有两个不同的解。因此在模$n$的意义上总共有$4$个解。
同时，这也就解决了上一章的一个遗留问题。
\end{example}
}

\frame{\frametitle{Example.} 
	\begin{example}[Example 7.]
		To compute $12^{12} \bmod 143 $ \pause
		解：
		\begin{itemize}
		\item 记$n = 143$，$n = p \cdot q = 13 \cdot 11$；
		
		\item 求：$12^{12} \bmod 13$ 和 $12^{12} \bmod 11 $；
		
		\item 所以，$12^{12} \bmod 143  = 1$.
		\end{itemize}
	\end{example}
	
}


\frame{\frametitle{Little thought.} 
	\begin{remark}
		A practical application: if we have many computations to perform on $x \in \mathbb{Z}_n^*$
		(e.g. RSA signing and decryption), we can convert $x$ to $(a, b) \in \mathbb{Z}_p^* \times \mathbb{Z}_q^*$
		and do all the computations on $a$ and $b$
		instead before converting back. 
		
		This is often cheaper because for many algorithms, doubling the size of the input more than doubles the running time.
	\end{remark}
}


\frame{\frametitle{Homework.} 
	\begin{block}{Homework.}
		1. Use CRT to solve:
		\begin{align*}
		x &\equiv 8 \pmod{11} \\
		x &\equiv 13 \pmod{19}
		\end{align*}
		
		2. Use CRT to solve the system of congruences:
			\begin{align*}
			x &\equiv 1 \pmod{5}\\
			x &\equiv 2 \pmod{7}\\
			x &\equiv 3 \pmod{9}\\
			x &\equiv 4 \pmod{11}
		\end{align*}
		% ans = 1731
		
		3. Write a program (in C or Python) to solve CRT.	
	\end{block}

}

\frame{\frametitle{Homework.} 
	\begin{block}{Homework.}
		4. Complete the proof that  it is an isomorphism from $\mathbb{Z}_n^*$ to $\mathbb{Z}_p^* \times \mathbb{Z}_q^*$.
		
		5. Let $p = 5$ and $q = 7$, $n = p \cdot q$. Please explicitly give the correspondence between $\mathbb{Z}_n^*$ and $\mathbb{Z}_p^* \times \mathbb{Z}_q^*$. Hint: Programming is permitted.
	\end{block}
	
}

\frame{\frametitle{Homework.} 
	\begin{block}{Homework.}
		6. 求解$x^{113} \equiv 15 \pmod{221}$[提示：这是一道似曾相识的习题，然而中国剩余定理能保证你顺利完成目标。]
		
	\end{block}
	
}
\end{document}
